\uuid{vyqe}
\exo7id{7179}
\titre{exo7 7179}
\auteur{megy}
\organisation{exo7}
\datecreate{2017-07-26}
\isIndication{true}
\isCorrection{true}
\chapitre{Propriétés de R}
\sousChapitre{Autre}
\cpge{mpsi-mp2i.calcul-algebrique-trigonometrie.c, pcsi.calcul-algebrique-trigonometrie.c, ptsi.calcul-algebrique-trigonometrie.c, mpsi-mp2i.fonctions-valeurs-reelles-complexes.c, pcsi.techniques-calcul-fonctions.c, ptsi.techniques-calcul-fonctions.c, bcpst1.nombres-reels, bcpst1.fonctions-reelles-usuelles.c}
\module{Analyse}
\niveau{L1}
\difficulte{}

\contenu{
\question{
Résoudre sur $\R$ l'équation
\[ 2^x + x^2 = 2-\frac{1}{2^x}.\]
}
\indication{Utiliser l'inégalité arithmético-géométrique.}
\reponse{
L'équation est équivalente à 
\[ 2^x +\frac{1}{2^x} = 2-x^2.\]
Or, par inégalité arithmético-géométrique, on a 
\[ 2^x +\frac{1}{2^x} \geq 2\sqrt{2^x\cdot \frac{1}{2^x}} \geq 2,\]
avec égalité ssi $2^x = \frac{1}{2^x}$ c'est-à-dire ssi $x=0$.

D'autre part, $2-x^2\geq 2$ avec égalité ssi $x=0$ là aussi.

On en déduit que l'équation admet bien une solution, unique, égale à $0$.
}
}
